The simulator does not select arbitrary activity labels at a decision point.
The BPMN engine converts the current case marking into executable transition
candidates. Each candidate contains a BPMN transition identifier, its visible
activity label, and, where required, the silent-transition path used to reach
it. Every branching decision is therefore filtered against the candidates
enabled under the current marking.

The runtime policy applies the predictive component first. If it returns no
BPMN-valid decision, the simulator uses the probabilistic component and finally
samples uniformly from the remaining enabled candidates. The offline
evaluation measures the predictive classifier separately, while the simulator
uses this composite chain to remain executable for unseen or unsupported
states.

The predictive component is a Random-Forest next-activity classifier
\cite{breiman2001randomforests,pedregosa2011scikit}. Its training data are
generated by replaying the event log on the BPMN model. An observation is
included only when the current event can be synchronized and fired, the
resulting marking enables several possible successors, and the observed next
activity belongs to that candidate set.

Rather than fitting a fully independent classifier for each gateway, we
train one parameter-shared Random Forest on replay-derived decision
observations. The current and previous activities, case-prefix information,
temporal and resource features, non-constant runtime attributes, and the
BPMN-enabled candidate set represent the decision context. Inference
therefore remains decision-point-specific even though model parameters are
shared. Predictions are restricted to transitions enabled under the current
marking, allowing the model to share statistical strength across sparsely
observed states while remaining BPMN-valid.

The probabilistic component implements the Basic branching method. It first
uses a state-conditioned successor distribution based on the current activity,
previous activity, and coarse repetition information. If the state has
insufficient training support, it falls back to the global successor
distribution of the current activity. If no supported successor remains after
BPMN filtering, the final fallback samples uniformly from the enabled
candidates.

The implementation distinguishes event-log activity labels from BPMN
transition identifiers, since a label may be ambiguous or temporarily
non-executable. Synchronization coverage, next-activity prediction, and exact
transition execution are therefore evaluated separately. Detailed results are
reported in Sec.~\ref{ssub:branching-eval} and
App.~\ref{app:joao_branching_model_selection}.